% ─────────────────────────────────────────────────────────────────────────────
% CHAPTER 6 — CONCLUSION AND FUTURE SCOPE
% ─────────────────────────────────────────────────────────────────────────────

\chapter{Conclusion and Future Scope}

The work presented in this report addresses the critical gap in maritime 
surveillance where existing systems rely either on \gls{ais} telemetry alone, 
which cannot detect non-cooperative vessels, or on isolated vision-based 
pipelines, which lack the navigational context to support collision 
avoidance. The proposed MARVIS system demonstrates that these two 
complementary data sources can be fused in real time within a unified 
deep-learning pipeline to deliver simultaneous vessel detection, identity 
tracking, trajectory prediction, and collision risk alerting. This chapter 
summarises the conclusions drawn from the design, implementation, and 
evaluation of the system, identifies the current limitations and their 
implications for practical deployment, and outlines directions for future 
extension.

% ─────────────────────────────────────────────────────────────────────────────
\section{Conclusion}
% ─────────────────────────────────────────────────────────────────────────────

Maritime surveillance in congested waterways such as the Singapore Strait 
demands a system capable of detecting all vessels regardless of whether 
they carry \gls{ais} transponders, assigning persistent identities across 
frames despite occlusions and viewpoint changes, predicting where each 
vessel is heading over the next several seconds, and raising timely 
collision risk alerts to support operator decision-making. Each of these 
requirements was established as a specific objective of this project: 
(1)~real-time vessel detection from drone or \gls{cctv} video streams, 
(2)~persistent multi-object tracking with appearance-based 
\gls{reid}, (3)~geo-referencing of detected vessels to \gls{gps} 
coordinates and fusion with live \gls{ais} telemetry, (4)~\gls{lstm}-based trajectory 
prediction with kinematic fallback for nascent tracks, and 
(5)~\gls{cpa}-based collision risk detection with a graded four-level alert 
scheme, all integrated into a dual-mode web dashboard accessible in 
real time.

Each objective was realised through a modular six-stage software pipeline 
implemented entirely in Python. The \gls{yolov8} anchor-free detection backbone 
was fine-tuned on the Roboflow maritime dataset and deployed with a 
confidence threshold of 0.4 to classify seven vessel types from 
frame-level video input. The \gls{deepocsort} multi-object tracker, combining 
a per-track Kalman filter with \gls{osnet} \gls{reid} embeddings and a 
cosine similarity threshold of 0.5, maintained persistent vessel 
identities across frames under occlusion and crossing conditions. 
The \texttt{GPSConverter} module transformed pixel centroids to 
geographic coordinates using a \gls{fov} linear approximation anchored 
to the Singapore Strait camera position, and the \texttt{CVMatcher} 
module performed nearest-neighbour \gls{ais} fusion using the Haversine 
great-circle distance formula within a 5.5\,km matching radius, enriching 
camera detections with \gls{mmsi}, vessel name, \gls{sog}, and course 
over ground from the live \gls{ais} stream. The \gls{seqtwo} \gls{lstm} predictor, 
comprising a bidirectional encoder, soft attention layer, and 
autoregressive LSTMCell decoder, was trained on Singapore Maritime Dataset 
trajectory sequences using a composite loss penalising positional error, 
velocity smoothness, and acceleration consistency, with a three-tier 
inference fallback ensuring that a trajectory estimate is produced for 
every tracked vessel regardless of track age. The \gls{cpa}-based collision 
detector evaluated all vessel pairs at each frame and raised graded 
alerts with a 30-frame cooldown to avoid alert fatigue.

The experimental evaluation on \gls{smd} test sequences confirms that the 
proposed system meets all five design objectives. The \gls{yolov8} detection 
module achieves an overall mAP\textsubscript{50} of 
\textit{[TBD from friend's results]}, demonstrating effective 
generalisation across vessel classes, scales, and sea-state conditions. 
The \gls{deepocsort} tracker achieves a \gls{mota} of 
\textit{[TBD]}\% with only \textit{[TBD]} identity switches across all 
test sequences, confirming the effectiveness of the combined 
Kalman-filter--\gls{reid} association in dense maritime scenes. The full \gls{lstm} 
trajectory predictor achieves an \gls{ade} of \textit{[TBD]} pixels over the 
20-step prediction horizon, representing a \textit{[TBD]}\% improvement 
over the kinematic linear regression baseline, with the advantage most 
pronounced for vessels executing turning manoeuvres and lane changes 
that linear extrapolation cannot anticipate. The \gls{cpa} collision detector 
achieves a precision of \textit{[TBD]} at the critical risk level 
($<50$\,m \gls{cpa}), providing actionable advance warning with sufficient 
lead time for operator intervention. Collectively, these results 
establish that a fully software-based, open-source maritime surveillance 
system integrating vision, \gls{ais} telemetry, deep learning prediction, and 
collision avoidance can be realised on commodity hardware and deployed 
with a single startup command.

% ─────────────────────────────────────────────────────────────────────────────
\section{Future Scope}
% ─────────────────────────────────────────────────────────────────────────────

While the proposed system achieves its stated objectives, several 
constraints and limitations in the current implementation present 
opportunities for extension in future work.

The geo-referencing module uses a \gls{fov} linear approximation 
that assumes a flat, undistorted projection centred on a fixed camera 
position. This is sufficient for small surveillance areas of up to 
2\,km radius but introduces increasing positional error at larger scales 
or under non-nadir viewing angles. Future work should replace this 
approximation with a full homographic transformation using four or more 
ground-control points with known \gls{gps} coordinates, which has already been 
prototyped in the \texttt{geo\_mapper.py} module. Automatic camera 
calibration from \gls{gps} metadata embedded in drone SRT files would further 
eliminate the need for manual anchor point configuration.

The current \gls{cpa} collision detection module operates on a 
constant-velocity motion model, using only the instantaneous heading 
and speed of each vessel. A significant extension would be to replace 
the constant-velocity \gls{cpa} prediction with the \gls{lstm}-predicted future 
trajectories, computing \gls{cpa} over the full 20-step predicted path rather 
than a linear velocity extrapolation. This would allow the system to 
anticipate collision risks arising from vessels that are currently 
diverging but whose predicted turning manoeuvres place them on a 
converging course in the near future.

The \gls{lstm} model was trained exclusively on the Singapore Maritime Dataset, 
which is captured in a specific geographic and traffic context. 
Generalisation to other maritime environments with different vessel 
density, lane structure, or vessel class distributions may require 
domain adaptation or transfer learning from a more diverse multi-region 
training corpus. Additionally, the current model architecture generates 
a single deterministic predicted trajectory per vessel; replacing this 
with a generative model such as a Conditional Variational Autoencoder 
(CVAE) or a Generative Adversarial Network (GAN) would enable the system 
to represent the multi-modal nature of vessel future paths, providing 
a distribution over possible trajectories rather than a single estimate.

The system currently runs on a conventional \gls{cpu}/\gls{gpu} workstation. 
For deployment on drone-mounted edge hardware, such as the \gls{nvidia} Jetson 
Nano or Jetson Orin, the inference pipeline would require quantisation 
of the \gls{yolov8} model to \gls{int8} precision and optimisation of the \gls{lstm} 
inference using TensorRT, both of which are supported by the existing 
\gls{pytorch} and Ultralytics toolchain. Edge deployment would eliminate the 
need for video transmission to a remote server, reducing latency and 
enabling autonomous surveillance in communication-denied environments. 
Finally, extension of the \gls{ais} ingestion module to support additional 
data sources, including radar returns and satellite \gls{ais}, would 
significantly increase the system's coverage and reliability in offshore 
open-sea surveillance scenarios.

% ─────────────────────────────────────────────────────────────────────────────
\section{Learning Outcomes of the Project}
% ─────────────────────────────────────────────────────────────────────────────

\begin{itemize}

    \item Gained practical experience in designing and implementing a 
    real-time multi-stage deep learning pipeline that integrates computer 
    vision, object tracking, sequence modelling, and geospatial data 
    fusion within a single production-grade software system.

    \item Developed a thorough understanding of the \gls{yolov8} anchor-free 
    object detection architecture, including transfer learning and 
    fine-tuning on a domain-specific maritime dataset, confidence and 
    \gls{iou} threshold selection, and the trade-offs between model variant 
    size (nano, medium, large) and inference throughput.

    \item Acquired proficiency in multi-object tracking algorithms, 
    specifically the \gls{deepocsort} framework, encompassing Kalman filter 
    state estimation, Hungarian assignment for detection-to-track 
    association, appearance-based \gls{reid} using \gls{osnet} embeddings, 
    and evaluation using CLEAR MOT metrics (\gls{mota}, \gls{motp}, \gls{ids}).

    \item Designed, trained, and evaluated a \gls{seqtwo} \gls{lstm} 
    model with bidirectional encoding and soft attention for non-linear 
    trajectory prediction, gaining experience in composite loss function 
    design, early stopping, the \gls{ade}/\gls{fde} evaluation protocol, and the 
    practical three-tier inference fallback strategy for varying track 
    history lengths.

    \item Learned to implement sensor fusion between heterogeneous data 
    sources --- video-based \gls{gps} estimates and live \gls{ais} telemetry --- using 
    the Haversine great-circle distance formula for nearest-neighbour 
    spatial matching, and understood the operational significance of 
    dark vessel detection for maritime domain awareness.

    \item Understood and implemented the \gls{cpa} 
    algorithm used in professional maritime ARPA systems, applying it 
    within a software collision detection module with graded risk 
    classification and alert cooldown logic, and gained awareness of 
    COLREGS collision regulations that govern maritime navigation.

    \item Gained experience in full-stack system integration, including 
    the development of a FastAPI and Flask backend with WebSocket 
    real-time data broadcast, a Folium-based interactive geospatial map 
    frontend, a Plotly-based analytics dashboard, and a Streamlit-based 
    MARVIS operator interface, and understood the trade-offs between 
    polling and push-based data delivery for real-time applications.

    \item Developed skills in collaborative software development using 
    Git version control, including branch management, conflict resolution, 
    pull request workflows, and the discipline of maintaining a stable 
    main branch while active development proceeds on feature branches.

\end{itemize}